\chapter{Verification plan: procedures and results}
\label{ch:verifplan}

Chapter~\ref{ch:verif} states the verification strategy and the headline
results. This chapter is the working document behind it: for each block-level
bench, what is being proved, the procedure step by step, and the pass criterion.
It is written so that a reader can follow a procedure without the source open,
and so that an engineer picking the work up can re-derive why each step is
there.

\section{How to read a procedure}

Every procedure below has the same shape.

\begin{table}[H]
\centering
\begin{tabular}{@{}lL{0.66\linewidth}@{}}
\toprule
\textbf{Objective} & The property being proved, stated as something that could
be false. \\
\textbf{Setup} & What is instantiated and how it is driven. \\
\textbf{Procedure} & Numbered steps. Each step is either a stimulus (``write
X'', ``raise line Y'') or an observation with its expected value. \\
\textbf{Pass criterion} & What must hold for the section to pass. \\
\bottomrule
\end{tabular}
\end{table}

Three conventions are used throughout and are worth stating once, because they
are what makes the difference between a procedure and a demonstration.

\paragraph{Every step is self-checking.} Observations are compared in the bench
with the \verb|===| operator, so an X never passes as a match. A mismatch prints
a \texttt{FAIL:} line and increments an error counter; the run ends on a
pass/fail banner that names the bench.

\paragraph{Negative properties get a witness.} ``No transaction was issued'' and
``the interrupt was never taken'' cannot be observed directly, so the benches
count what \emph{did} happen --- AXI handshakes, acknowledge pulses, handler
entries --- and require the count to be zero. A count is falsifiable; an absence
of output is not.

\paragraph{Every negative is paired with a positive.} A test that shows nothing
happened is worthless if nothing could have happened. Wherever a procedure
proves that something was correctly suppressed, the next step removes the
suppressing condition and requires the thing to happen. That is what
distinguishes masking from a dead path, and it is called out explicitly in each
procedure below.

\section{Regression composition}

\begin{table}[H]
\centering
\caption{The twelve regression runs, from one compile}
\label{tab:reg12}
\begin{tabular}{@{}clL{0.40\linewidth}rl@{}}
\toprule
\textbf{\#} & \textbf{Bench} & \textbf{Scope} & \textbf{Checks} & \textbf{Result} \\
\midrule
1  & \file{tb\_cpu\_axi}    & System, default latencies (CPI check active here) & 90 & Pass \\
2  & \file{tb\_cpu\_axi}    & System, high fixed AXI latencies & 90 & Pass \\
3  & \file{tb\_cpu\_axi}    & System, random backpressure, seeds 101/202 & 90 & Pass \\
4  & \file{tb\_cpu\_axi}    & System, random backpressure, seeds 777/888 & 90 & Pass \\
5  & \file{tb\_dual\_core}  & Two cores, shared memory, 2:1 arbiter & --- & Pass \\
6  & \file{tb\_pic}         & PIC feature bench & 139 & Pass \\
7  & \file{tb\_pic\_reset}  & PIC reset values, X-freedom, async reset & 350 & Pass \\
8  & \file{tb\_pic\_ro}     & PIC read-only registers and reserved bits & 478 & Pass \\
9  & \file{tb\_pic\_status} & PIC \reg{SRCx\_STATUS}, field by field & 424 & Pass \\
10 & \file{tb\_mtimer\_regs}& Machine-timer registers and access rules & 224 & Pass \\
11 & \file{tb\_csr\_ro}     & CSR read-only / WARL / tied-off / absent & 282 & Pass \\
12 & \file{tb\_traps}       & One directed test per trap cause & 111 & Pass \\
\midrule
   & \textbf{Total}         & & \textbf{2\,278} & \textbf{12/12} \\
\bottomrule
\end{tabular}
\end{table}

Compile and elaboration report \texttt{Errors: 0, Warnings: 0} for every run.
The banners are deliberately distinct per bench, so the log is validated by
counting them: a bench that dies early prints no banner \emph{and} no
\texttt{FAIL:} line, which a ``no failures'' check alone would accept.

% ===========================================================================
\section{PIC reset verification}
\label{sec:vp-picreset}

\textbf{Bench:} \file{debug/hdl/tb\_pic\_reset.v} --- 350 checks.

\paragraph{Objective.} Prove that every flop in the interrupt controller is
reset by \sig{rst\_n\_i}; that every mapped register returns its documented reset
value on its first access; that nothing leaves reset holding X, including state
with no register view; that the reset is asynchronous and dominant; and that the
block leaves reset disarmed.

\paragraph{Why it is a separate bench.} A reset defect is a first-cycle defect.
The feature bench spends its opening hundred cycles configuring the block, so by
the time it observes anything the reset values are gone. This bench observes
before anything is written.

\paragraph{Setup.} The PIC alone, with its AXI4-Lite port driven by the shared
master tasks, its 16 source lines and its claim/end-of-interrupt pins driven
directly by the bench. Clock 10\,ns; reset is under the bench's control rather
than the shared generator's, because the procedure asserts it a second time.

\subsection*{Procedure}

\begin{enumerate}[leftmargin=*,itemsep=4pt]

\item[\textbf{1}] \textbf{Bus and CPU outputs while reset is asserted.}
\begin{enumerate}[label=1.\arabic*,leftmargin=12mm,itemsep=1pt]
  \item Hold \sig{rst\_n\_i} = 0 for four clock cycles.
  \item Sample the slave outputs. \sig{BVALID} and \sig{RVALID} must be 0
        (IHI0022E A3.1.2); \sig{BRESP}, \sig{RRESP} and \sig{RDATA} must contain
        no X and must read \texttt{OKAY}, \texttt{OKAY} and
        \texttt{0x0000\_0000}.
  \item Sample the CPU side: \sig{cpu\_irq} = 0, \sig{cpu\_irq\_vec} = 0, both
        X-free.
\end{enumerate}

\item[\textbf{2}] \textbf{Every mapped register reads its reset value.}
\begin{enumerate}[label=2.\arabic*,leftmargin=12mm,itemsep=1pt]
  \item Release \sig{rst\_n\_i} between clock edges.
  \item Read all 16 \reg{SRCx\_CONFIG}, all 16 \reg{SRCx\_SW\_TRIG}, all 16
        \reg{SRCx\_STATUS} and the 11 global registers --- 59 registers --- and
        compare each against Table~\ref{tab:picreset}.
\end{enumerate}
All 16 instances of each per-source register are read, not a sample. A register
that resets correctly at index 0 and not at index 11 is exactly what a spot
check misses.

\item[\textbf{3}] \textbf{CPU-side outputs after release} are inactive and
X-free.

\item[\textbf{4}] \textbf{State with no register view.} Probe all 16
\sig{stack\_id} and \sig{stack\_key} entries through hierarchical references and
require them X-free; require \sig{depth} = 0. These are the flops a ``the valid
bit covers it'' argument would leave unreset, and they are read by
\reg{NEST\_STATUS.TOP\_ID} and \reg{ACTIVE\_VEC.ID}.

\item[\textbf{5}] \textbf{Asynchronous reset from a fully loaded state.}
\begin{enumerate}[label=5.\arabic*,leftmargin=12mm,itemsep=1pt]
  \item Write every global register to a value different from its reset value:
        \reg{BAND\_CONFIG} = \texttt{0x27}, \reg{NEST\_MAX} = 4,
        \reg{ESCALATION\_CFG} = \texttt{0x110}, \reg{INT\_ENABLE} =
        \texttt{0xFFFF}.
  \item Write all 16 \reg{SRCx\_CONFIG} and set all 16 software channels.
  \item Raise two source lines and claim two nesting levels.
  \item \emph{Confirm the precondition:} nesting depth must be non-zero and
        \reg{BAND\_CONFIG} must read back \texttt{0x27}. Without this step the
        section could pass on a block that was never loaded.
  \item Assert \sig{rst\_n\_i} = 0 \emph{off a clock edge}. \sig{cpu\_irq} must
        drop before the next edge.
  \item Hold three cycles, release between edges.
  \item Repeat the full 59-register read-back of step 2, and the stack probe of
        step 4.
\end{enumerate}

\item[\textbf{6}] \textbf{The block is disarmed out of reset.}
\begin{enumerate}[label=6.\arabic*,leftmargin=12mm,itemsep=1pt]
  \item With no configuration written at all, hold source 4 high for eight
        cycles.
  \item \sig{cpu\_irq} must stay low and \reg{SRC4\_STATUS.PEND} must stay 0,
        because \reg{INT\_ENABLE} reset to 0.
  \item \emph{The paired positive:} write \reg{INT\_ENABLE} = \texttt{0x0010}
        with the same line still high. \sig{cpu\_irq} must now assert. This is
        what proves step 6.2 observed masking rather than a dead source line.
\end{enumerate}
\end{enumerate}

\paragraph{Pass criterion.} 350 checks, zero failures. Step 5 must produce the
same 59 values as step 2.

\paragraph{Result.} Pass. Run against the previous revision of the RTL, step
1.2 fails on \sig{BRESP} and \sig{RRESP}, and step 4 fails on all 32
nesting-stack entries. That is the evidence that the procedure detects the
defect it was written for.

% ===========================================================================
\section{PIC read-only and reserved-bit verification}
\label{sec:vp-picro}

\textbf{Bench:} \file{debug/hdl/tb\_pic\_ro.v} --- 478 checks.

\paragraph{Objective.} Prove that every read-only object in the register map is
read-only, and that every reserved bit position is hardwired to zero.

\paragraph{The two promises.} ``Read-only'' means two separate things, and a
test that checks only one of them is not a test:

\begin{enumerate}[leftmargin=*,itemsep=2pt]
\item the access is \textbf{rejected} --- the write answers SLVERR, so software
finds out instead of believing it succeeded;
\item the state is \textbf{unchanged} --- the register still holds what the
hardware put there.
\end{enumerate}

The second promise is vacuous unless the register holds a non-trivial value at
the moment of the attempt. Writing 0 to a register that already reads 0 and then
observing 0 proves nothing. Every read-only register in this bench is therefore
driven to a known non-zero value by real hardware activity first.

\subsection*{Procedure}

\begin{enumerate}[leftmargin=*,itemsep=4pt]

\item[\textbf{1}] \textbf{\reg{SRC0..15\_STATUS} --- 16 read-only registers.}
\begin{enumerate}[label=1.\arabic*,leftmargin=12mm,itemsep=1pt]
  \item Enable all 16 sources.
  \item Give every source the maximum deadline (\texttt{0xFFFF}), so
        \sig{DDL\_TIMER} runs --- making the register hardware-driven and
        non-zero --- without ever expiring inside the run.
  \item Raise all 16 lines; let the timers advance.
  \item For each index: read, write all-ones, require SLVERR, read again, and
        require the six state bits \sig{PEND}, \sig{ACTIVE}, \sig{ESC},
        \sig{SPUR} and \sig{EFF\_BAND} to be identical.
  \item For \sig{DDL\_TIMER}, which is a live counter and cannot be checked by
        ``unchanged'', require the stronger property: after the rejected write
        it must have \emph{advanced}, and it must not equal the written value.
  \item Repeat step 1.4 writing \texttt{0x0000\_0000}: a rejection may not
        depend on the data.
\end{enumerate}

\item[\textbf{2}] \textbf{\reg{NEST\_STATUS}.} Claim one interrupt so
\sig{DEPTH} and \sig{TOP\_ID} are both non-zero, then apply the read-only
sequence.

\item[\textbf{3}] \textbf{\reg{ACTIVE\_VEC}.} With the interrupt still in
service, so \sig{VALID} and \sig{ID} are set, apply the same sequence.

\item[\textbf{4}] \textbf{Unmapped offsets.} Every word from \texttt{0xE0} to
\texttt{0xFC} must answer SLVERR on \emph{both} a read and a write. This is what
turns a stray pointer into a precise access-fault trap in the CPU instead of a
silent read of zero.

\item[\textbf{5}] \textbf{Reserved bits.} For each register, write all-ones and
require the reserved positions to read back 0 while the live positions keep the
value:
\begin{center}\small
\begin{tabular}{@{}ll@{}}
\reg{SRCx\_CONFIG} (all 16) & reserved [15:8] and [3] \\
\reg{BAND\_CONFIG}          & reserved [31:8] \\
\reg{NEST\_MAX}             & reserved [31:5] \\
\reg{ESCALATION\_CFG}       & reserved [31:9], [7:5] and [3:2] \\
\reg{INT\_ENABLE}           & reserved [31:16] \\
\reg{SPURIOUS\_LOG}         & reserved [31:16] \\
\reg{INT\_STATUS}           & reserved [31:3] \\
\reg{SRCx\_SW\_TRIG} (all 16) & only bit [0] is readable state \\
\end{tabular}
\end{center}

\item[\textbf{6}] \textbf{\reg{SRCx\_SW\_TRIG.KEY} is write-only.} Set a channel
with the real key, read the register back, and require the key field to read 0.
If the key read back, software could discover it by reading the register, which
defeats the purpose of having one.

\item[\textbf{7}] \textbf{Write-1-to-clear semantics.} Produce a spurious
claim, so that bit 6 of \reg{SPURIOUS\_LOG} and the \sig{SPUR} bit of
\reg{INT\_STATUS} are both set by hardware. Then: write 0 to both, and nothing
may change; write 1 to the \emph{wrong} bit, and the set bit must survive;
write 1 to the right bit, and it must clear. This is the
other half of being read-only: a bit that software cannot \emph{set}. A read,
modify and write back that puts a 0 into the register must leave a set bit
alone, or every handler that touches the register loses events it never meant
to acknowledge.

\item[\textbf{8}] \textbf{\reg{NEST\_MAX} is WARL.} Write 0 (stores 1), 17
(stores 16), 31 (stores 16), 5 (stores 5 exactly). The clamp is not cosmetic: a
depth limit of 0 would mask every offer for ever.

\item[\textbf{9}] \textbf{Narrow writes.} Attempt a byte write to each of the
four lanes of a read-only register, plus a write with no lanes enabled: all five
must SLVERR and change nothing. \emph{The paired positive:} the same byte-lane
write to a \emph{writable} register must succeed, proving the rejection was
decoded from the address and not from the strobe.
\end{enumerate}

\paragraph{Pass criterion.} 478 checks, zero failures.

\begin{deviation}
\textbf{Finding.} This procedure found a real deviation between the delivered
RTL and the register specification: the reserved bits of \reg{SRCx\_CONFIG}
([15:8] and [3]) and of \reg{ESCALATION\_CFG} ([7:5] and [3:2]) were stored and
read back rather than hardwired to zero. Software could therefore write a 1 into
a bit the specification called reserved and read it back, acquiring a dependency
that a later revision defining that bit would break. The RTL was corrected to
mask them on write, matching the documented map; the design is otherwise
unchanged and the feature bench required no modification.
\end{deviation}

% ===========================================================================
\section{PIC status-register verification}
\label{sec:vp-picstatus}

\textbf{Bench:} \file{debug/hdl/tb\_pic\_status.v} --- 424 checks.

\paragraph{Objective.} \reg{SRCx\_STATUS} is the only window software has into
what a source is doing. Section~\ref{sec:vp-picro} proves software cannot write
it; this proves the hardware puts the right thing in it.

\paragraph{Method note.} Every field is checked in \emph{both} directions: it
must come up when its condition holds and go away when the condition stops
holding. A status bit stuck at 1 is as useless as one stuck at 0, and only the
two-sided check finds the first.

\begin{table}[H]
\centering
\caption{\reg{SRCx\_STATUS} verification, field by field}
\label{tab:vpstatus}
\begin{tabular}{@{}L{0.15\linewidth}L{0.37\linewidth}L{0.32\linewidth}@{}}
\toprule
\textbf{Field} & \textbf{Procedure} & \textbf{Key point} \\
\midrule
\sig{PEND} [0] &
  1. line high with \reg{INT\_ENABLE} = 0 $\rightarrow$ must stay 0.
  2. enable the source $\rightarrow$ must become 1.
  3. drop the line $\rightarrow$ must return to 0.
  4. a neighbouring source must not have moved. &
  The enable term is the part worth proving: software must be able to tell an
  ignored peripheral from a quiet one. \\
\midrule
\sig{ACTIVE} [1] &
  1. raise a level source $\rightarrow$ \sig{PEND}=1, \sig{ACTIVE}=0.
  2. claim $\rightarrow$ both 1, because the line is still high.
  3. clear the source in the handler $\rightarrow$ \sig{PEND}=0,
     \sig{ACTIVE}=1.
  4. end of interrupt $\rightarrow$ both 0. &
  \sig{PEND} and \sig{ACTIVE} are independent, and step 2 is where a design that
  conflated them would fail. \\
\midrule
\sig{DDL\_TIMER} [31:16] &
  1. idle $\rightarrow$ 0. 2. pending $\rightarrow$ advancing, measured twice.
  3. claimed $\rightarrow$ cleared, and stays 0 for the whole handler.
  4. idle again $\rightarrow$ 0. 5. a source with no deadline configured holds
  it at 0 while genuinely pending. &
  The counter measures how long the source \emph{waited for service}, so it must
  stop the moment service starts. \\
\midrule
\sig{ESC} [2] and \sig{EFF\_BAND} [5:4] &
  1. band 3, deadline 8, jump-to-band-0 policy. 2. hold the request past the
  deadline. 3. \sig{ESC}=1 \emph{and} \sig{EFF\_BAND}=0. 4. \reg{SRCx\_CONFIG}
  must be untouched. 5. after service both revert. &
  Checked together on purpose: a design that set the flag and forgot to move the
  band would silently lose the whole feature. \\
\midrule
\sig{EFF\_BAND}, bump mode &
  Watch the effective band on \emph{every} clock edge, count the transitions
  and flag any that moved more than one band. Require exactly three
  transitions, none skipping, ending at band 0; then require it to saturate
  rather than wrap. &
  A register read costs several cycles, so polling over AXI cannot tell a
  correct $3\rightarrow2\rightarrow1\rightarrow0$ walk from a single
  $3\rightarrow0$ jump. This is why the check is a per-cycle monitor. \\
\midrule
\sig{SPUR} [3] &
  \emph{Negative first:} a source still asserted at the claim and cleared later
  inside the handler must \textbf{not} be flagged, and must not touch
  \reg{SPURIOUS\_LOG}. \emph{Then positive:} drop the line between the offer and
  the claim $\rightarrow$ \sig{SPUR}=1, sticky log bit set,
  \reg{INT\_STATUS.SPUR} set, nesting depth still 1 (the claim is accounted, so
  ack and eoi stay balanced), and on return \sig{SPUR} clears while the sticky
  log does not. &
  The distinction is \emph{when} the request disappears, not how fast. The
  negative case is the normal behaviour of every peripheral in this system, so
  getting it wrong would flag every interrupt. \\
\bottomrule
\end{tabular}
\end{table}

\paragraph{Life-cycle walk.} The bench closes with one source taken through its
whole life, reading the register at every transition, so the fields are shown as
a sequence rather than as six independent facts.

\begin{figure}[H]
\centering
\begin{tikzpicture}[node distance=17mm]
  \node[st] (idle) {IDLE};
  \node[st,right=of idle] (pend) {PENDING};
  \node[st,right=of pend] (svc)  {IN\\SERVICE};
  \node[st,right=of svc]  (done) {SERVICED};
  \draw[flow] (idle) -- node[lbl,above]{line high} (pend);
  \draw[flow] (pend) -- node[lbl,above]{claim} (svc);
  \draw[flow] (svc)  -- node[lbl,above]{handler} node[lbl,below]{clears source} (done);
  \draw[flow] (done.south) .. controls +(0,-11mm) and +(0,-11mm) ..
     node[lbl,below]{end of interrupt} (idle.south);
  \node[below=6mm of idle,font=\scriptsize,align=center] {all 0\\timer 0};
  \node[below=6mm of pend,font=\scriptsize,align=center] {PEND=1\\timer running};
  \node[below=6mm of svc,font=\scriptsize,align=center]  {PEND=1 ACTIVE=1\\timer cleared};
  \node[below=6mm of done,font=\scriptsize,align=center] {ACTIVE=1\\timer 0};
\end{tikzpicture}
\caption{\reg{SRCx\_STATUS} life cycle. The bench reads the register at each of
the four states and after the return, and compares the whole six-bit state
field, so a bit that is right in one state and wrong in the next is caught.}
\label{fig:statuslife}
\end{figure}

\paragraph{Pass criterion.} 424 checks, zero failures.

% ===========================================================================
\section{Machine-timer register verification}
\label{sec:vp-mtimer}

\textbf{Bench:} \file{debug/hdl/tb\_mtimer\_regs.v} --- 224 checks.

\paragraph{Objective.} Verify the timer as a block in its own right rather than
only through the system bench, for two reasons: its reset values are the unusual
ones in the delivery (\reg{mtimecmp} resets to all-ones, so the timer is born
disarmed), and its only defence against an out-of-range access is the shared
slave's address decode, which nothing else here exercises.

\subsection*{Procedure}

\begin{enumerate}[leftmargin=*,itemsep=3pt]
\item[\textbf{1}] \textbf{During reset:} \sig{BVALID}/\sig{RVALID} low,
\sig{BRESP}/\sig{RRESP}/\sig{RDATA} X-free, \sig{irq} low.

\item[\textbf{2}] \textbf{Reset values.} \reg{MTIME\_HI} = 0;
\reg{MTIME\_LO} within a small window of 0, because it is already counting;
\reg{MTIMECMP\_LO} and \reg{MTIMECMP\_HI} both all-ones; \sig{irq} low. Then let
it run 200 cycles: a disarmed timer must never fire.

\item[\textbf{3}] \textbf{Rate.} Measure the counter delta between two
back-to-back reads, then between two reads separated by 50 extra clock cycles.
The difference between the two deltas must be \emph{exactly} 50. Taking the
difference cancels the fixed AXI overhead, which turns an approximate check into
an exact one.

\item[\textbf{4}] \textbf{Round-trip} on each register independently, with a
cross-check after each write that the other three are untouched. The
\reg{MTIME\_LO} case is checked against a window opening at the written value,
because the counter resumes immediately.

\item[\textbf{5}] \textbf{Byte strobes.} Write all-ones one lane at a time and
require exactly that lane to change; then a write with no lanes enabled, which
must answer OKAY and change nothing.

\item[\textbf{6}] \textbf{Unmapped offsets.} All 60 words from \texttt{0x10} to
\texttt{0xFC} must SLVERR on read and on write. \emph{The paired positive:} the
four mapped words must still answer OKAY, so step 6 is a decode result and not a
blanket rejection.

\item[\textbf{7}] \textbf{Interrupt behaviour.}
\begin{enumerate}[label=7.\arabic*,leftmargin=12mm,itemsep=1pt]
  \item Disarm, zero the counter, confirm \sig{irq} low.
  \item Arm in the documented order --- \reg{MTIMECMP\_LO} first, then
        \reg{MTIMECMP\_HI}. After the low half alone, \sig{irq} must still be
        low, because the high half is still all-ones and the compare cannot be
        satisfied. This is the step that proves the arming order is safe.
  \item Write the high half; \sig{irq} must rise when \reg{mtime} reaches
        \reg{mtimecmp}.
  \item \sig{irq} is a level: it must hold high while the compare holds.
  \item Move \reg{mtimecmp} forward, as a handler would; \sig{irq} must clear.
        There is no status register to write --- the compare \emph{is} the
        status.
  \item \emph{The paired positive for 7.2:} set \reg{mtimecmp} to a time already
        past. It must fire at once, proving the hardware compares the full 64
        bits and that step 7.2 was not simply a timer that never fires.
\end{enumerate}

\item[\textbf{8}] \textbf{Asynchronous reset} from a configured, armed, firing
state, asserted off a clock edge. \sig{irq} must drop before the next edge, and
all four registers must return to their reset values --- in particular
\reg{mtimecmp} back to all-ones, leaving the timer disarmed again.
\end{enumerate}

\paragraph{Pass criterion.} 224 checks, zero failures.

% ===========================================================================
\section{CSR access-rule verification}
\label{sec:vp-csr}

\textbf{Bench:} \file{debug/hdl/tb\_csr\_ro.v} --- 282 checks.

\paragraph{Objective.} The CSR file has four access classes that are easy to
confuse, and confusing them is the classic CSR defect. This bench separates them
and checks each on its own terms.

\begin{table}[H]
\centering
\caption{CSR access classes and the rule each must obey}
\label{tab:csrclass}
\begin{tabular}{@{}lL{0.28\linewidth}L{0.34\linewidth}@{}}
\toprule
\textbf{Class} & \textbf{Registers} & \textbf{Required behaviour on a write} \\
\midrule
Read-only & \reg{mip}, \reg{mvendorid}, \reg{marchid}, \reg{mimpid},
  \reg{mhartid} & \textbf{Must trap} (illegal instruction) and leave the value
  alone. The trap is what lets software discover the mistake. \\
WARL & \reg{misa}; \reg{mepc}[1:0]; the \reg{mstatus} WPRI fields;
  \reg{mie}[15:0] & \textbf{Must not trap.} The write is accepted and only the
  legal part is kept. A WARL field that traps breaks conforming software. \\
Tied off & \reg{mhpmcounter8..31}, their upper halves, \reg{mhpmevent3..31} ---
  77 addresses & \textbf{Must not trap}, must read 0, must swallow the write.
  They exist and are hardwired, which is not the same as absent. \\
Absent & everything else, including \reg{mcountinhibit} & \textbf{Must trap} on
  a read \emph{and} on a write. \\
\bottomrule
\end{tabular}
\end{table}

\subsection*{Procedure}

\begin{enumerate}[leftmargin=*,itemsep=3pt]
\item[\textbf{1}] \textbf{Reset value of every implemented CSR}, against
Table~\ref{tab:csrreset}. \reg{mstatus} is the one that is not zero, because MPP
is hardwired to \texttt{11}.

\item[\textbf{2}] \textbf{Read-only CSRs.} \reg{mip} is first driven non-zero by
asserting two interrupt lines, so ``unchanged'' means something. Each of the
five is then attacked with all three CSR operations --- CSRRW, CSRRS and
CSRRC --- each of which must raise illegal, and the value re-read afterwards.

\item[\textbf{3}] \textbf{Read-only \emph{fields} inside writable CSRs.} These
are the bits a register-level test misses: the CSR accepts the write, nothing
traps, but part of the value must not survive.
\begin{center}\small
\begin{tabular}{@{}lll@{}}
\toprule
\textbf{CSR} & \textbf{Written} & \textbf{Must read back} \\
\midrule
\reg{mstatus} & \texttt{0xFFFF\_FFFF} & \texttt{0x0000\_1888} (MIE, MPIE, MPP=11) \\
\reg{mie}     & \texttt{0xFFFF\_FFFF} & \texttt{0xFFFF\_0000} \\
\reg{mie}     & \texttt{0x0000\_FFFF} & \texttt{0x0000\_0000} (no live bit touched) \\
\reg{mepc}    & \texttt{0xFFFF\_FFFF} & \texttt{0xFFFF\_FFFC} \\
\bottomrule
\end{tabular}
\end{center}
\emph{The paired positive:} \reg{mtvec} and \reg{mscratch} are written with
arbitrary 32-bit patterns and must keep every bit, showing that the rows above
masked fields rather than dropping writes.

\item[\textbf{4}] \textbf{\reg{misa} is WARL, not read-only.} All three
operations must be accepted without trapping, and the value must stay
\texttt{0x4000\_0100}. This is the distinction the bench exists to make: a
design that made \reg{misa} read-only would trap here and break conforming
startup code that probes the extension set by writing to it.

\item[\textbf{5}] \textbf{Tied-off registers.} All 77 addresses read 0, accept a
write of all-ones without trapping, and still read 0 afterwards.

\item[\textbf{6}] \textbf{Absent CSRs} raise illegal on read and on write:
\reg{mcountinhibit} (deliberately not implemented --- the specification makes it
optional and defines the absent behaviour as ``all counters run'', which is what
this design does), \reg{sstatus}, \reg{pmpcfg0}, the user-mode cycle counter,
and the addresses \texttt{0x000} and \texttt{0xFFF}. The \emph{boundaries} of
each tied-off range are then checked one address either side, which is where an
off-by-one in a range comparison would show.

\item[\textbf{7}] \textbf{The pure-read exception.} CSRRS or CSRRC with
\texttt{rs1} = x0 is defined as a read with no write side effect, so it must not
trap even on a read-only CSR.

\item[\textbf{8}] \textbf{Implemented counters are writable} on both halves,
independently --- the other side of step 5: tied off is not the same as
read-only.

\item[\textbf{9}] \textbf{Asynchronous reset} from a fully written state returns
every CSR to its reset value, while \reg{mhartid} and \reg{misa} are unaffected
because they are hardwired rather than stored.
\end{enumerate}

\paragraph{Pass criterion.} 282 checks, zero failures.

% ===========================================================================
\section{Trap-cause verification}
\label{sec:vp-traps}

\textbf{Bench:} \file{debug/hdl/tb\_traps.v} --- 111 checks.

\paragraph{Objective.} Give every trap cause its own isolated test, with the
four architectural results checked individually: \reg{mcause}, \reg{mepc},
\reg{mtval}, and the handler address actually entered.

\paragraph{Why, given the system bench already covers traps.} The system bench
proves every cause is reachable and counts the entries exactly. That is the
right test for ``the trap mechanism works end to end''. It is the wrong test for
``cause 6 records the right \reg{mtval}'', because a single failing cause is one
line in a hundred, and because the causes are interleaved with everything else
the program does. Isolation also makes the negative half checkable: for every
cause that must not commit its instruction or must not issue a bus transaction,
that is asserted in the same case instead of inferred.

\paragraph{Setup.} \reg{cpu\_top} against two behavioural AXI4-Lite memories:
instruction memory at \texttt{0x0000} (256 words, so \texttt{0x8000} is
unmapped) and data memory at \texttt{0x2000} (256 words, so \texttt{0x4000} is
unmapped). The interrupt pins are driven directly by the bench rather than
through the PIC, so the CPU's trap logic is isolated from the controller's.

Each case is one reset, one small program written directly into the instruction
memory word by word, and one trap. The program always begins by setting
\reg{mtvec}; the handler is a self-loop, which is what makes ``which handler was
entered'' observable. Architectural results are read from the CSR file through
hierarchical references, so the check is exact without adding a debug port to
the RTL.

\begin{figure}[H]
\centering
\begin{tikzpicture}[node distance=4.5mm]
  \node[blkf,minimum width=96mm] (s1) {\textbf{1.} assert reset, write the case
    program into instruction memory};
  \node[blkf,below=of s1,minimum width=96mm] (s2) {\textbf{2.} release reset
    between clock edges; the program sets \reg{mtvec}};
  \node[blkf,below=of s2,minimum width=96mm] (s3) {\textbf{3.} the case
    instruction executes; exactly one thing can go wrong and it is the thing
    under test};
  \node[blkf,below=of s3,minimum width=96mm] (s4) {\textbf{4.} wait for
    \sig{cpu\_in\_trap}; check \reg{mcause}, \reg{mepc}, \reg{mtval}};
  \node[blkf,below=of s4,minimum width=96mm] (s5) {\textbf{5.} let the pipeline
    settle on the handler self-loop; check the entry address};
  \node[blk,below=of s5,minimum width=96mm] (s6) {\textbf{6.} check the negative
    property for this cause: no writeback, or no bus transaction};
  \draw[flow] (s1) -- (s2); \draw[flow] (s2) -- (s3); \draw[flow] (s3) -- (s4);
  \draw[flow] (s4) -- (s5); \draw[flow] (s5) -- (s6);
\end{tikzpicture}
\caption{The procedure applied to every case in Table~\ref{tab:trapcases}.}
\label{fig:trapproc}
\end{figure}

{\small
\begin{longtable}{@{}cL{0.25\linewidth}lllL{0.19\linewidth}@{}}
\caption{Trap-cause test cases. \texttt{H} is the handler base (\texttt{0x0000\_0080}).}
\label{tab:trapcases}\\
\toprule
\textbf{Cause} & \textbf{Stimulus} & \reg{mepc} & \reg{mtval} & \textbf{Entry} & \textbf{Also checked} \\
\midrule
\endfirsthead
\toprule
\textbf{Cause} & \textbf{Stimulus} & \reg{mepc} & \reg{mtval} & \textbf{Entry} & \textbf{Also checked} \\
\midrule
\endhead
\bottomrule
\endfoot
\bottomrule
\endlastfoot
0  & \texttt{jal x0,+6} at \texttt{0x08}; target \texttt{0x0E} has bit~1 set
   & \texttt{0x08} & \texttt{0x0E} & \texttt{H} & no bus traffic \\
1  & \texttt{jal} to \texttt{0x8000}, outside the instruction memory
   & \texttt{0x8000} & \texttt{0x8000} & \texttt{H} & \reg{mepc} is the address
     that could not be fetched, not the jump \\
2  & the word \texttt{0xFFFF\_FFFF} (opcode \texttt{1111111})
   & \texttt{0x08} & \texttt{0xFFFF\_FFFF} & \texttt{H} & no writeback,
     no bus traffic \\
3  & \texttt{EBREAK}
   & \texttt{0x08} & \texttt{0x08} & \texttt{H} & \reg{mepc} points \emph{at}
     the EBREAK, so a debugger can resume it \\
11 & \texttt{ECALL}
   & \texttt{0x08} & \texttt{0} & \texttt{H} & no faulting address exists, so
     \reg{mtval} is 0 \\
4  & \texttt{lw x5,1(x1)} with \texttt{x1}=\texttt{0x2000}, a \emph{mapped}
     address & \texttt{0x0C} & \texttt{0x2001} & \texttt{H} & zero bus
     transactions; \texttt{x5} untouched \\
4  & \texttt{lh x5,1(x1)} --- the halfword rule is bit~0 only
   & \texttt{0x0C} & \texttt{0x2001} & \texttt{H} & zero bus transactions \\
6  & \texttt{sw x0,2(x1)}, \texttt{x1}=\texttt{0x2000}
   & \texttt{0x0C} & \texttt{0x2002} & \texttt{H} & zero bus transactions ---
     the dangerous case, since a store that reached the bus would have corrupted
     memory \\
6  & \texttt{sh x0,1(x1)}
   & \texttt{0x0C} & \texttt{0x2001} & \texttt{H} & zero bus transactions \\
5  & \texttt{lw x5,0(x1)} with \texttt{x1}=\texttt{0x4000}, aligned but unmapped
   & \texttt{0x0C} & \texttt{0x4000} & \texttt{H} & exactly one read
     transaction; \texttt{x5} untouched \\
7  & \texttt{sw x0,0(x1)} with \texttt{x1}=\texttt{0x4000}
   & \texttt{0x0C} & \texttt{0x4000} & \texttt{H} & exactly one write
     transaction \\
\midrule
\multicolumn{6}{@{}l@{}}{\emph{Cause priority}} \\
4  & \texttt{lw x5,3(x1)} with \texttt{x1}=\texttt{0x4000}: the address is
     \emph{both} misaligned and unmapped & \texttt{0x0C} & \texttt{0x4003} &
     \texttt{H} & cause 4 must win over cause 5, with zero bus transactions \\
\midrule
\multicolumn{6}{@{}l@{}}{\emph{Machine external interrupts}} \\
16 & \reg{mstatus.MIE} = 1 and \reg{mie}[16] = 1; request on vector 0
   & boundary & \texttt{0} & \texttt{H} & \reg{mcause} =
     \texttt{0x8000\_0010}; acknowledge pulsed once \\
19 & \reg{mie}[19] only; vector 3 & boundary & \texttt{0} & \texttt{H} &
     \reg{mcause} = \texttt{0x8000\_0013} \\
31 & \reg{mie}[31] only; vector 15 & boundary & \texttt{0} & \texttt{H} &
     \reg{mcause} = \texttt{0x8000\_001F} \\
\end{longtable}}

For the interrupt cases \reg{mepc} is required to be a 4-aligned address inside
the program body rather than a fixed number, because it is the instruction to
\emph{resume} at and which instruction that is depends on when the request
arrived. Everything else about the interrupt is exact.

\paragraph{Cause priority, and why it needs its own case.} Alignment is checked
before the access is issued, so an address that is both misaligned and unmapped
must raise cause 4 and issue nothing. Without this case, an implementation that
issued first and trapped on the response afterwards would still pass every other
row in the table.

\begin{figure}[H]
\centering
\begin{tikzpicture}[node distance=3.6mm]
  \node[blk,minimum width=104mm] (a) {interrupt pending, enabled, and at an
     instruction boundary?};
  \node[blkf,below=of a,minimum width=104mm] (b) {instruction access fault
     (cause 1) --- the fetch already failed};
  \node[blkf,below=of b,minimum width=104mm] (c) {illegal instruction
     (cause 2)};
  \node[blkf,below=of c,minimum width=104mm] (d) {breakpoint (cause 3)};
  \node[blkf,below=of d,minimum width=104mm] (e) {environment call (cause 11)};
  \node[blkf,below=of e,minimum width=104mm] (f) {instruction address
     misaligned (cause 0)};
  \node[blkf,below=of f,minimum width=104mm] (g) {load / store address
     misaligned (causes 4, 6) --- \emph{checked before any access is issued}};
  \node[blkf,below=of g,minimum width=104mm] (h) {load / store access fault
     (causes 5, 7) --- only knowable from the bus response};
  \node[blk,below=of h,minimum width=104mm] (i) {no trap: the instruction
     commits};
  \foreach \x/\y in {a/b,b/c,c/d,d/e,e/f,f/g,g/h,h/i} \draw[flow] (\x) -- (\y);
\end{tikzpicture}
\caption{Trap priority as implemented. The interrupt is tested first and at an
instruction boundary, so a preempted instruction reruns cleanly after
\texttt{MRET}. The two shaded groups near the bottom are the pair the
cause-priority case separates.}
\label{fig:trappri}
\end{figure}

\subsection*{Negative cases}

Two negatives are checked, each immediately followed by the positive that proves
it was masking and not a dead path.

\begin{enumerate}[leftmargin=*,itemsep=3pt]
\item \textbf{A vector masked in \reg{mie} is never taken.} \reg{mstatus.MIE}=1
and \reg{mie} enables cause 16 only; a request is raised on vector 5 and held
for 300 cycles. \sig{cpu\_in\_trap} must stay low, the acknowledge must never
pulse, and \reg{mcause} must never be written. \emph{Then} the same request is
switched to vector 0, and must be taken at once with
\reg{mcause}=\texttt{0x8000\_0010}. Without this pairing, an implementation that
ignored the vector entirely and took any request would still pass the three
positive interrupt cases.

\item \textbf{A request with \reg{mstatus.MIE}=0 is never taken.} \reg{mie}
enables cause 16, but the global enable is left at its reset value of 0. The
enabled vector is held for 300 cycles with no trap and no acknowledge.
\end{enumerate}

\subsection*{Vectored \reg{mtvec}}

Both halves of the vectored rule are checked. A design that sent
\emph{everything} to BASE\,+\,4$\cdot$cause would pass a check that looked
only at interrupts.

\begin{enumerate}[leftmargin=*,itemsep=2pt]
\item \reg{mtvec} = \texttt{H}\,|\,1; take interrupt cause 19. The handler entry
must be \texttt{H}\,+\,76.
\item \reg{mtvec} = \texttt{H}\,|\,1; take \texttt{ECALL}, cause 11. The handler
entry must be \texttt{H} --- vectoring applies to interrupts only.
\end{enumerate}

\paragraph{Pass criterion.} 111 checks, zero failures; every case must reach its
trap inside the 2\,000-cycle limit.
